% Figure: adsorption breakthrough curve, outlet concentration vs time.
% S-shaped rise; breakthrough at the design limit, saturation at the feed value.
\begin{figure}[htbp]
\centering
\begin{tikzpicture}
\begin{axis}[
  width=11cm, height=6.5cm,
  xlabel={time (or throughput volume)},
  ylabel={outlet/feed concentration $c_{\text{out}}/c_{\text{in}}$},
  domain=0:10, samples=120,
  xmin=0, xmax=10, ymin=0, ymax=1.1,
  xtick=\empty,
  legend pos=north west,
  legend cell align=left,
  grid=both, grid style={enggray!18},
  every axis plot/.append style={very thick},
]
% logistic breakthrough curve centred at t=6
\addplot[engblue] {1/(1+exp(-2.2*(x-6)))};
\addlegendentry{breakthrough curve}
% design breakthrough limit
\draw[engred, dashed, thick] (axis cs:0,0.05) -- (axis cs:10,0.05);
\node[font=\scriptsize, engred, anchor=west] at (axis cs:0.2,0.16) {design limit};
% breakthrough time marker
\draw[enggray, dashed] (axis cs:4.6,0) -- (axis cs:4.6,0.05);
\node[font=\scriptsize, anchor=north] at (axis cs:4.6,-0.02) {$t_b$};
% saturation marker
\draw[enggray, dashed] (axis cs:8.2,0) -- (axis cs:8.2,0.95);
\node[font=\scriptsize, anchor=north] at (axis cs:8.2,-0.02) {$t_{\text{sat}}$};
\end{axis}
\end{tikzpicture}
\caption[Adsorption breakthrough curve]{The breakthrough curve of a
fixed-bed adsorber. While fresh sorbent remains ahead of the mass-transfer
zone, the outlet is clean; once the zone reaches the bed exit the outlet
concentration climbs and the bed approaches saturation at the feed value.
The bed is taken offline at the breakthrough time $t_b$, when the outlet
crosses the design limit. A sharp mass-transfer zone gives a steep curve
and high sorbent utilisation; a broad zone wastes capacity and risks early
breakthrough.}
\label{fig:vol04ch07:breakthrough}
\end{figure}
